%%%%%%%%%%%%%%%%%%%%%%%%%%%%%%%%%%%%%%%%%%%%%%%%%%%%%%%%%%%%%%%%%%%%%%%%%%%%%%%
%% Two Roads to Quantum Safety: A Hybrid E91-QKD + ML-DSA Architecture
%% with CHSH/QBER Eavesdropping Detection and Threshold Co-Optimisation
%%
%% IEEE format submission, v2.3 (LaTeX, conference format), June 2026 — revised per reviewer critique
%% Authors: Vinit Datta Rane, Shuchona Malek Orthi, Mehrdad S. Sharbaf
%%
%% Compile with pdfLaTeX (Overleaf default). Bibliography uses IEEEtran.bst.
%%%%%%%%%%%%%%%%%%%%%%%%%%%%%%%%%%%%%%%%%%%%%%%%%%%%%%%%%%%%%%%%%%%%%%%%%%%%%%%

\documentclass[conference]{IEEEtran}
\IEEEoverridecommandlockouts

% ---------- Packages ----------
\usepackage{cite}
\usepackage{amsmath,amssymb,amsfonts}
\usepackage{algorithmic}
\usepackage{graphicx}
\usepackage{textcomp}
\usepackage{xcolor}
\usepackage{array}
\usepackage{booktabs}
\usepackage{url}
\usepackage{hyperref}
\usepackage{orcidlink}
\usepackage{siunitx}

% IEEE Access encourages hyperref but with specific settings to avoid
% conflicts with the IEEEtran class.
\hypersetup{
    colorlinks=true,
    linkcolor=blue,
    citecolor=blue,
    urlcolor=blue,
    pdftitle={Two Roads to Quantum Safety: A Hybrid E91-QKD + ML-DSA Architecture},
    pdfauthor={Vinit Datta Rane, Shuchona Malek Orthi, Mehrdad S. Sharbaf}
}

% ---------- Document ----------
\begin{document}

\title{Two Roads to Quantum Safety: A Hybrid E91-QKD + ML-DSA Architecture with CHSH/QBER Eavesdropping Detection and Threshold Co-Optimisation
\thanks{Corresponding author: Vinit Datta Rane (vrane1@toromail.csudh.edu). This work received no external funding.}}

\author{
\IEEEauthorblockN{Vinit Datta Rane\,\orcidlink{0009-0004-9740-8009}}
\IEEEauthorblockA{\textit{MS Cybersecurity Program} \\
\textit{Calif. State Univ., Dominguez Hills}\\
Carson, CA, USA \\
vrane1@toromail.csudh.edu}
\and
\IEEEauthorblockN{Shuchona Malek Orthi\,\orcidlink{0009-0007-5397-4561}}
\IEEEauthorblockA{\textit{MS IT Project Management Program} \\
\textit{Westcliff University}\\
Irvine, CA, USA \\
s.orthi.339@westcliff.edu}
\and
\IEEEauthorblockN{Mehrdad S. Sharbaf, Ph.D.}
\IEEEauthorblockA{\textit{Computer Science Department} \\
\textit{Calif. State Univ., Dominguez Hills}\\
Carson, CA, USA \\
msharbaf@csudh.edu}
}

% IEEE Access running header removed per revision.

\maketitle

% ============================================================
% ABSTRACT (IEEE Access: single unstructured paragraph, <=250 words)
% ============================================================
\begin{abstract}
A key question in quantum-safe networking is how to combine Quantum Key Distribution (QKD) and Post-Quantum Cryptography (PQC) so that each covers the other's principal failure mode. This paper specifies a hybrid architecture pairing the entanglement-based E91 QKD protocol with ML-DSA (FIPS 204, Level 3) authentication and evaluates it through discrete-event simulation under realistic SPDC parameters (100\,MHz pair rate, 10\% coupling, 85\% detector efficiency, 0.2\,dB/km fiber). Four contributions are reported. First, we identify the KEM/signature design distinction: key encapsulation primitives (ML-KEM, HQC) cannot authenticate basis announcements, so ML-DSA is the correct primitive for the classical channel. Second, we present a six-step protocol with a qualified composability argument that follows Bindel \emph{et~al.}'s hybrid combiner theorem by analogy. Third, we characterise CHSH inequality testing and QBER monitoring as partially complementary detectors with different operational sensitivity regions, with QBER dominating at low intercept fractions and CHSH activating analytically at $p \geq 0.29$. Fourth, we present finite-key analysis identifying a $27.4$\,km hybrid/PQC crossover, and a 2D threshold co-optimisation showing that CHSH=$2.66$ and QBER=$2.33\%$ raise simulated detection probability at $p = 0.05$ from $64.5\%$ to $75.6\%$ ($+11.1$ percentage points absolute, $17.2\%$ relative) while holding the false-positive rate below $0.5\%$. A 100\,000-trial empirical validation at the optimised operating point yields a false-positive rate of $0.392\%$ (95\% Wilson CI $[0.355\%, 0.433\%]$), confirming the analytical $0.40\%$ prediction. The hybrid is practical for metropolitan deployments up to $50$\,km; ML-DSA overhead is below $0.25\%$ per protocol round.
\end{abstract}

% IEEE Access: 4-6 Index Terms drawn from IEEE Thesaurus, first letter capitalised, comma separated, no period after last term
\begin{IEEEkeywords}
Quantum key distribution, Post-quantum cryptography, Cryptographic protocols, Authentication, Network security
\end{IEEEkeywords}

% ============================================================
\section{Introduction}
% ============================================================
\IEEEPARstart{T}{he} digital world runs on encryption. Every online purchase, every private message, and every login passes through algorithms designed to keep data secret from anyone but the intended receiver. For decades, this worked. The mathematical problems behind RSA and Elliptic Curve Cryptography were hard enough that no practical computer could break them in a human lifetime. That assumption is now under threat. Quantum computers, once a theoretical curiosity, have moved into real hardware. Although still too small to crack real-world keys, steady progress in qubit count and error correction suggests that a cryptographically relevant quantum computer may appear within ten to fifteen years~\cite{mosca2018}.

The security community has prepared two distinct answers. The first is Quantum Key Distribution (QKD), which uses the laws of quantum physics to let two parties generate a shared random secret key~\cite{bennett1984,ekert1991}. Any attempt to eavesdrop on the quantum channel disturbs the transmitted photons, providing physics-based, computational-assumption-free eavesdropping detection with finite statistical confidence bounded by the verification block size. The second answer is Post-Quantum Cryptography (PQC): new mathematical problems designed to remain hard even against quantum computers, deployable as a software upgrade on existing infrastructure~\cite{nist2022}.

The two are sometimes framed as competitors. That framing is wrong. Prior work on hybrid cryptographic schemes has demonstrated that combining primitives from different security foundations yields improved resilience against failure in any single component~\cite{bindel2019,stebila2010}. QKD provides eavesdropping detection but requires an authenticated classical channel; PQC can provide that authentication. In turn, QKD provides PQC with physical randomness for key refresh.

\textbf{Contributions.} This paper makes four contributions:

\begin{itemize}
\item \textbf{KEM/signature design clarification.} We articulate that KEMs (ML-KEM, HQC) cannot authenticate arbitrary message content and that ML-DSA is the correct primitive for QKD classical channel authentication.
\item \textbf{Reference protocol with qualified composability argument.} We specify a complete six-step E91+ML-DSA protocol with an informal security argument following Bindel \emph{et~al.}~\cite{bindel2019} by analogy.
\item \textbf{Simulation-based detection characterisation.} We report key rates, $|S|$, QBER, and detection probabilities with 95\% confidence intervals, including analytical derivation of the CHSH/QBER crossover threshold at $p \approx 0.29$.
\item \textbf{Finite-key analysis, crossover, and threshold co-optimisation.} We identify the $27.4$\,km hybrid/PQC crossover, present a 2D threshold co-optimisation grid, and empirically validate the optimised operating point at $N = 10^5$ trials.
\end{itemize}

The paper is structured as follows. Section~\ref{sec:background} provides background. Section~\ref{sec:related} surveys related work. Section~\ref{sec:architecture} describes the hybrid architecture. Section~\ref{sec:simulation} presents simulation results. Section~\ref{sec:discussion} discusses limitations and deployment. Section~\ref{sec:conclusion} concludes.

% ============================================================
\section{Background}
\label{sec:background}
% ============================================================

\subsection{The Quantum Threat to Classical Cryptography}
Most internet security today relies on RSA and Elliptic Curve Cryptography. Both are trapdoor functions: easy forward, hard to reverse without a secret key. Shor's algorithm~\cite{shor1997} factors integers in polynomial time on a sufficiently large quantum computer, breaking both. Grover's algorithm~\cite{grover1996} accelerates symmetric key search, reducing AES-128's effective security to 64 bits and motivating migration to AES-256. ``Harvest now, decrypt later'' attacks make the timeline urgency greater than the direct attack timeline suggests~\cite{mosca2018}.

\subsection{Quantum Key Distribution}
QKD uses individual photons to transmit a random key without relying on computational hardness~\cite{gisin2002}. BB84~\cite{bennett1984} uses single-photon polarisation states with two bases. The E91 protocol~\cite{ekert1991} uses entangled photon pairs measured in three bases per party. Eavesdropping is detected via the CHSH inequality~\cite{chsh1969}:
\begin{equation}
|S| = |E(a,b) - E(a,b') + E(a',b) + E(a',b')| \leq 2 \label{eq:chsh}
\end{equation}
where $E(a,b)$ is the expectation of the product of outcomes in bases $a$ and $b$. Maximally entangled photons yield $|S| = 2\sqrt{2} \approx 2.828$, violating the classical bound. Any eavesdropper who intercepts and resends photons reduces the degree of entanglement, lowering $|S|$ toward the classical bound of $2$, while introducing errors measurable as elevated QBER.

Throughout this paper we model Eve as an individual attacker performing random-basis intercept-resend on a fraction $p$ of the photons. Under this canonical individual-attack strategy~\cite{fuchs1997}, the expected CHSH expectation value of the intercepted fraction averages to $\sqrt{2}$ over Eve's basis choice, which leads to the linear interpolation~\eqref{eq:S_of_p} derived in Section~\ref{sec:methodology}. More sophisticated collective and coherent attacks~\cite{scarani2009}, as well as hardware-side attacks~\cite{lydersen2010,makarov2009}, exist but are out of scope for the quantitative simulation.

It is important to distinguish the CHSH abort threshold from the QKD security threshold. The security threshold is the maximum QBER at which privacy amplification can still extract a secure key (typically $\sim 11\%$ for E91~\cite{scarani2009}). The abort threshold is set more conservatively to detect eavesdropping early. The security of QKD does not depend on Eve's computational power~\cite{scarani2009}, though real devices have imperfections (photon loss, dark counts, source imperfections, misalignment)~\cite{scarani2009,xu2020,treiber2009}.

This paper focuses on discrete-variable (DV) QKD. Continuous-variable (CV) QKD~\cite{pirandola2020} offers easier telecom integration; adapting the hybrid architecture to CV-QKD is left as future work.

\subsection{Post-Quantum Cryptography and the KEM/Signature Distinction}
NIST's post-quantum standardisation~\cite{nist2022} produced FIPS~203 (ML-KEM)~\cite{fips203} for key encapsulation and FIPS~204 (ML-DSA)~\cite{fips204} for digital signatures. HQC~\cite{hqc2022} was selected as an additional KEM finalist with code-based hardness.

A critical distinction exists between these two classes. Key encapsulation mechanisms (KEMs) such as ML-KEM and HQC establish a shared symmetric key but do not authenticate arbitrary message content. Digital signature schemes such as ML-DSA produce unforgeable signatures over arbitrary data. For QKD classical channel authentication, which requires signing basis announcements, CHSH data, and error correction messages, a signature scheme is the correct primitive. Using a KEM for this purpose is a design error. Our protocol uses ML-DSA explicitly to avoid it.

% ============================================================
\section{Related Work}
\label{sec:related}
% ============================================================

\subsection{Foundational QKD}
Bennett--Brassard's BB84~\cite{bennett1984} and Ekert's E91~\cite{ekert1991} established the foundations of QKD. Security proofs against arbitrary attacks were developed under ideal device assumptions~\cite{scarani2009}. Practical QKD advanced via decoy-state protocols~\cite{lo2005}, CV-QKD~\cite{pirandola2020}, and satellite-based entanglement distribution over 1{,}200\,km~\cite{liao2017}. Comprehensive reviews appear in Pirandola \emph{et~al.}~\cite{pirandola2020}, Xu \emph{et~al.}~\cite{xu2020}, and Gisin \emph{et~al.}~\cite{gisin2002}.

\subsection{NIST PQC Standardisation}
The NIST process~\cite{nist2022} produced FIPS~203 (ML-KEM)~\cite{fips203} and FIPS~204 (ML-DSA)~\cite{fips204} in 2024. HQC~\cite{hqc2022} was selected as an additional code-based KEM. Standardised, vetted algorithms make software integration of PQC straightforward.

\subsection{Hybrid QKD+PQC: Prior Work and Positioning}
Stebila, Mosca, and L\"utkenhaus argued foundationally that QKD without authenticated classical channels is trivially breakable by man-in-the-middle, and that PQC provides the natural authentication complement~\cite{stebila2010}. The ETSI QKD Application Interface~\cite{etsi2020} formalises this need without prescribing the primitive. Bindel \emph{et~al.}~\cite{bindel2019} provided formal UC proofs for hybrid KEM combiners.

More recently, Garms \emph{et~al.}~\cite{garms2024} delivered the first experimental demonstration of a hybrid QKD+PQC cryptosystem on a single FPGA. Zeng \emph{et~al.}~\cite{zeng2024} analysed multiple hybrid protocols across network topologies. Concurrently, Hsieh \emph{et~al.}~\cite{hsieh2025} proposed a hybrid scheme combining ML-KEM and ML-DSA with BB84 and E91.

Our work is positioned as follows. Garms \emph{et~al.}~\cite{garms2024} provided hardware implementation but did not analyse joint CHSH/QBER detection behaviour. Zeng \emph{et~al.}~\cite{zeng2024} provided formal network-level analysis but not simulation-based detection probability measurements. The concurrent E91+ML-DSA scheme of~\cite{hsieh2025} focuses on key exchange entropy without reporting detection behaviour. Our contribution is the explicit KEM/signature design principle, the simulation-based evaluation with realistic SPDC parameters and confidence intervals, and the quantitative characterisation of CHSH and QBER as complementary detectors with direct implications for abort threshold co-design. Mosca~\cite{mosca2018} has argued throughout that harvest-now-decrypt-later attacks make hybrid deployment urgently necessary today. Table~\ref{tab:positioning} summarises this positioning against the most directly comparable recent works.

\begin{table}[!t]
\caption{Positioning Against Recent Hybrid QKD+PQC Work (Authors' Own Compilation)}
\label{tab:positioning}
\centering
\small
\begin{tabular}{@{}lcccccc@{}}
\toprule
\textbf{Work} & \textbf{HW} & \textbf{CHSH} & \textbf{QBER} & \textbf{Finite-} & \textbf{Thresh.} & \textbf{Proto.} \\
 & & \textbf{anal.} & \textbf{anal.} & \textbf{key} & \textbf{opt.} & \textbf{spec.} \\
\midrule
Garms \emph{et~al.}~\cite{garms2024} & Yes & No & No & No & No & Partial \\
Zeng \emph{et~al.}~\cite{zeng2024} & No & No & No & Yes & No & No \\
Hsieh \emph{et~al.}~\cite{hsieh2025} & No & No & No & No & No & Yes \\
\textbf{This paper} & No & \textbf{Yes} & \textbf{Yes} & \textbf{Yes} & \textbf{Yes} & \textbf{Yes} \\
\bottomrule
\end{tabular}
\end{table}

% ============================================================
\section{Hybrid Architecture}
\label{sec:architecture}
% ============================================================

\subsection{Overview}
The hybrid system consists of two tightly coupled layers. The lower layer is the E91 QKD protocol~\cite{ekert1991}, generating a raw shared key via entangled photon measurements verified by both CHSH testing~\cite{chsh1969} and QBER monitoring. The upper layer is ML-DSA (FIPS~204, Security Level~3)~\cite{fips204}, authenticating all classical channel communication. HQC~\cite{hqc2022} may optionally encrypt the classical channel for confidentiality.

\begin{figure}[!t]
\centering
\includegraphics[width=\columnwidth]{figures/fig1_architecture.png}
\caption{Hybrid E91+ML-DSA system architecture. Alice's SPDC source emits entangled photon pairs; Alice retains photon$_A$ and transmits photon$_B$ to Bob over the quantum fiber channel. Eve's intercept-resend position is shown at the midpoint. All classical post-processing messages are ML-DSA signed (FIPS~204, Level~3). HQC may optionally encrypt the classical channel. (Authors' own work.)}
\label{fig:architecture}
\end{figure}

\subsection{Protocol Steps}
\textbf{Step 1: Mutual Authentication.} Alice and Bob hold long-term ML-DSA key pairs with public keys certified through a quantum-safe PKI. At the start of each round, each party signs a 256-bit nonce with its private key; mutual identity is established before any quantum transmission.

\textbf{Step 2: Quantum Transmission (E91).} An entangled photon source co-located with Alice emits photon pairs. Alice retains one photon per pair and transmits the other to Bob. Each party independently chooses measurement bases from $\{0^\circ, 45^\circ, 90^\circ\}$.

\textbf{Step 3: Basis Announcement and Sifting.} Over the ML-DSA--authenticated classical channel, Alice and Bob announce basis choices. They separate measurements into a verification subset (25\% of pulses, from non-compatible pairs) used for CHSH testing, and a sifted key subset used for key extraction.

\textbf{Step 4: Dual Detection Check.} Two independent abort conditions are evaluated. (a)~CHSH: compute $|S|$; abort if $|S| < 2.40$. (b)~QBER: estimate the quantum bit error rate; abort if QBER $> 2.5\%$. As quantified in Section~\ref{sec:simulation}, these checks have different sensitivity regions.

\textbf{Step 5: Error Correction and Privacy Amplification.} LDPC error correction~\cite{elkouss2011} on the sifted key, followed by privacy amplification via a SHA-256--based universal hash extractor.

\textbf{Step 6: Key Material Usage and Re-keying.} The final shared symmetric key drives AES-256-GCM. Periodically, the QKD-derived key material is fed through HKDF to seed a fresh ML-DSA key pair generated deterministically at both endpoints, then mutually authenticated over the current channel before switching. This turns long-term ML-DSA keys into dynamically refreshed session keys.

\emph{Note:} the re-keying mechanism is presented as an architectural design pattern rather than a standards-track specification; seed synchronisation, rollback handling, and stale-key revocation are deployment-specific concerns beyond the scope of this paper.

\begin{figure}[!t]
\centering
\includegraphics[width=\columnwidth]{figures/fig2_sequence.png}
\caption{Protocol sequence diagram for the hybrid E91+ML-DSA protocol. Coloured horizontal bands delineate each of the six steps. Navy arrows carry quantum photons; red arrows carry ML-DSA--signed authentication messages; green arrows carry authenticated classical post-processing messages; amber carries the periodic re-keying exchange. The ABORT branch at Step~4 is triggered by either the CHSH or QBER threshold. (Authors' own work.)}
\label{fig:sequence}
\end{figure}

\subsection{Security Argument}
We apply the hybrid combiner framework of Bindel \emph{et~al.}~\cite{bindel2019} by analogy. Their theorem proves, in the UC model, that a KEM combiner formed from a classically secure KEM and a post-quantum secure KEM is secure as long as at least one component is secure. We apply the combiner intuition to the hybrid authentication layer as an informal design argument (not a formal proof). Informally, the protocol is intended to remain operationally safe as long as at least one of the following conditions holds at the time of attack.

\emph{Quantum layer secure:} The E91 quantum channel provides physics-based, computational-assumption-free eavesdropping detection. Any interception disturbs entanglement, reducing $|S|$ and raising QBER. The detection mechanism does not depend on Eve's computational power; its statistical confidence at a given eavesdropping fraction is bounded by the verification block size.

\emph{Classical layer secure:} ML-DSA authentication is computationally secure. Eve cannot forge basis announcements without solving the Module-SIS lattice problem, which no known quantum algorithm solves efficiently~\cite{fips204}.

\emph{Important qualification:} The argument above is an informal security argument, not a formal proof in the Universal Composability model. A complete UC proof would require specifying ideal functionalities for both the QKD post-processing pipeline and the ML-DSA authentication scheme, and proving that the real protocol UC-realises the composed ideal functionality. Extending Bindel \emph{et~al.}'s theorem to QKD post-processing authentication is a non-trivial open problem identified for future work. The concise composable security framework of Lim \emph{et~al.}~\cite{lim2014} provides a suitable starting point for such a proof.

\subsection{Scope of Security Claims}
\label{sec:scope}
To delineate precisely what is and is not claimed, we state the following boundary explicitly:
\begin{itemize}
\item \textbf{No UC proof is provided.} We do not claim a proof in the Universal Composability model. A complete UC treatment would require specifying ideal functionalities for both the QKD post-processing pipeline and the ML-DSA authentication scheme and showing that the real protocol UC-realises their composition; we identify this as open work.
\item \textbf{The hybrid combiner analogy is conceptual.} Our use of the Bindel \emph{et~al.}~\cite{bindel2019} combiner theorem is an informal design argument by analogy, not a formal reduction. Extending it to QKD post-processing authentication is non-trivial; the composable frameworks of Renner~\cite{renner2008}, Lim \emph{et~al.}~\cite{lim2014}, and Portmann and Renner~\cite{portmann2022} provide suitable starting points.
\item \textbf{Stated security regime.} The protocol is argued secure under standard assumptions: ML-DSA (Module-SIS) hardness for the classical channel, and the canonical individual intercept-resend attack model for the quantum channel. Collective, coherent, and hardware-side attacks are out of scope for the quantitative analysis and are discussed qualitatively in Section~\ref{sec:limitations}.
\end{itemize}
Accordingly, this work should be interpreted as a systems-oriented hybrid architecture and simulation study rather than a formally verified cryptographic construction.

% ============================================================
\section{Simulation-Based Evaluation}
\label{sec:simulation}
% ============================================================

\subsection{Methodology and Exact Formulas}
\label{sec:methodology}
We built a discrete-event simulator in Python (seed~$=42$; publicly available at~\cite{orthi_rane_zenodo}) using parameters consistent with early telecom-fiber entanglement-based QKD experiments~\cite{poppe2004}. Three configurations are evaluated: Pure QKD (no classical channel authentication), Pure PQC (ML-KEM-768 key encapsulation), and Hybrid (E91+ML-DSA).

\textbf{Pure PQC baseline.} We model Pure PQC as repeated ML-KEM-768 (FIPS~203 Level~3)~\cite{fips203} key encapsulations followed by AES-256-GCM bulk encryption with periodic session-key refresh. Each ML-KEM-768 encapsulation produces 256~bits of shared key with a round-trip of $\sim 2.2$\,KB. The 122\,kbps Pure PQC entry in Table~\ref{tab:keyrates} represents the steady-state shared-secret refresh rate from repeated handshakes over a 1\,Gbps classical link, accounting for protocol framing, $\sim 3$\,ms metro-link round-trip latency, and key-derivation overhead. This is a key-refresh-rate metric, not a bulk-data-encryption throughput; the AES-256-GCM channel operates at full line rate independent of this metric.

\textbf{Coincidence rate.} Both photons must be detected:
\begin{equation}
R_{\mathrm{coin}} = R_{\mathrm{pair}} \cdot \eta_{\mathrm{coupling}} \cdot T(d)^2 \cdot \eta_{\mathrm{det}}^2 \label{eq:coin}
\end{equation}
where $R_{\mathrm{pair}} = 100$\,MHz, $\eta_{\mathrm{coupling}} = 0.10$, $T(d) = 10^{-0.2d/10}$ (fiber transmittance at distance $d$\,km), and $\eta_{\mathrm{det}} = 0.85$.

\textbf{Sifting.} E91 uses three bases per party with $P_{\mathrm{basis}} = 1/3$. We define the compatible-basis rate (key-extraction pairs plus CHSH-verification pairs) as
\begin{equation}
R_{\mathrm{compatible}} = R_{\mathrm{coin}} \cdot P_{\mathrm{basis}}. \label{eq:rcomp}
\end{equation}
Of the compatible-basis rounds, $V_{\mathrm{frac}} = 0.25$ is assigned to CHSH verification and $1 - V_{\mathrm{frac}} = 0.75$ to key generation.

\textbf{Quantum bit error rate:}
\begin{equation}
e_{\mathrm{bit}} = e_{\mathrm{misalign}} + e_{\mathrm{dark}} + p_{\mathrm{eve}}/3 \label{eq:qber}
\end{equation}
with $e_{\mathrm{misalign}} = 0.01$, $e_{\mathrm{dark}} \approx 10^{-5}$, and per-trial Gaussian noise $\sigma = 0.005$.

\textbf{CHSH parameter} under random-basis intercept-resend. Following the optimal individual-attack analysis of Fuchs \emph{et~al.}~\cite{fuchs1997}, an eavesdropper who measures an intercepted photon in a randomly chosen basis and resends destroys the entanglement on that fraction, so its contribution to the correlation averages to the separable value $\sqrt{2}$ rather than the Tsirelson value $2\sqrt{2}$. Linearly mixing the undisturbed fraction $(1-p_{\mathrm{eve}})$ with the intercepted fraction $p_{\mathrm{eve}}$ gives
\begin{equation}
|S(p)| = (1 - p_{\mathrm{eve}}) \cdot |S_{\mathrm{base}}| + p_{\mathrm{eve}} \cdot \sqrt{2} \label{eq:S_of_p}
\end{equation}
where $|S_{\mathrm{base}}| = 2\sqrt{2} \cdot (1 - e_{\mathrm{misalign}}) \approx 2.80$ and per-trial noise $\sigma = 0.04$. Setting $|S(p)| = 2.40$ yields the closed form $|S(p)| = 2.80 - 1.386\,p$, so the CHSH abort threshold is reached analytically at $p \approx 0.29$, consistent with the QBER--CHSH relationship under general attacks discussed by Fung \emph{et~al.}~\cite{fung2006}.

\textbf{Asymptotic secret key rate} (following Scarani \emph{et~al.}~\cite{scarani2009} and Lim \emph{et~al.}~\cite{lim2014}):
\begin{equation}
R_{\mathrm{secret}} = R_{\mathrm{compatible}} \cdot (1 - V_{\mathrm{frac}}) \cdot \max(0, 1 - f h(e_{\mathrm{bit}}) - h(e_{\mathrm{phase}})) \label{eq:rsecret}
\end{equation}
with $f = 1.16$ (LDPC inefficiency~\cite{elkouss2011}), $h(\cdot)$ the binary entropy function, and $e_{\mathrm{phase}} = e_{\mathrm{bit}}$ under the symmetric intercept-resend model. The $(1 - V_{\mathrm{frac}})$ factor selects the key-generation subset of $R_{\mathrm{compatible}}$, recovering the standard Scarani form $R_{\mathrm{secret}} = R_{\mathrm{keygen}} \cdot \max(0, 1 - f h(e_{\mathrm{bit}}) - h(e_{\mathrm{phase}}))$ by definition of $R_{\mathrm{keygen}} \equiv (1 - V_{\mathrm{frac}}) \cdot R_{\mathrm{compatible}}$.

\textbf{Finite-key length.} We report two estimators. The asymptotic secret fraction~\eqref{eq:rsecret} follows the Devetak--Winter rate~\cite{devetak2005}, which lower-bounds the extractable key by the difference between Alice--Bob mutual information and Eve's Holevo information. For finite blocks we apply the composable leftover-hash bound of Tomamichel \emph{et~al.}~\cite{tomamichel2012}, refined by Renner's smooth-entropy framework~\cite{renner2008}:
\begin{equation}
\ell_{\mathrm{rig}}(n) = n_{\mathrm{key}}\big[1 - f h(e_{\mathrm{bit}}) - h(e_{\mathrm{phase}})\big] - \Delta(n_{\mathrm{key}}) \label{eq:finite_rig}
\end{equation}
with the finite-size penalty
\begin{equation}
\Delta(n) = 7\sqrt{n\,\log_2(2/\bar{\varepsilon})} + 2\log_2\!\frac{1}{2\varepsilon_{\mathrm{PA}}} + \log_2\!\frac{2}{\varepsilon_{\mathrm{EC}}}, \label{eq:delta}
\end{equation}
where $n_{\mathrm{key}} = n_{\mathrm{raw}}(1 - V_{\mathrm{frac}})$, $\bar\varepsilon$ is the smooth-entropy parameter, and the leading $\sqrt{n}$ term carries the statistical-fluctuation cost of parameter estimation as in~\cite{tomamichel2012,tomamichel2017,hayashi2012}. For comparison and as a conservative engineering bound we also retain a simplified heuristic, $\ell_{\mathrm{heur}}(n) = n_{\mathrm{key}}\max(0, 1 - f h(e_{\mathrm{bit}}) - h(e_{\mathrm{bit}})) - 7\sqrt{n_{\mathrm{key}}}\log_2(2/\varepsilon_{\mathrm{sec}})$, with $\varepsilon_{\mathrm{sec}} = 10^{-10}$. The heuristic applies the full $\sqrt{n}$ penalty against the asymptotic key fraction, whereas the rigorous bound~\eqref{eq:finite_rig}--\eqref{eq:delta} separates reconciliation leakage from the smooth-entropy correction and applies the statistical fluctuation only to the phase-error estimate. As a result the heuristic substantially over-estimates the penalty: as reported in Section~\ref{sec:finitekey}, it yields $n_{\min} \approx 1.12\times10^5$ bits against the rigorous $n_{\min} \approx 9.9\times10^3$ bits. All $n_{\min}$ values reported in Section~\ref{sec:finitekey} are evaluated with the rigorous bound~\eqref{eq:finite_rig}; the heuristic value is quoted only as a conservative upper bound for comparison.

\textbf{Composable security budget.} The composable security parameter $\varepsilon_{\mathrm{sec}} = 10^{-10}$ decomposes by the union bound as $\varepsilon_{\mathrm{sec}} \leq \varepsilon_{\mathrm{PE}} + \varepsilon_{\mathrm{EC}} + \varepsilon_{\mathrm{PA}} + \varepsilon_{\mathrm{smooth}}$, with $\varepsilon_{\mathrm{PE}}$ bounding parameter-estimation failure, $\varepsilon_{\mathrm{EC}}$ bounding reconciliation failure, $\varepsilon_{\mathrm{PA}}$ bounding privacy-amplification leakage, and $\varepsilon_{\mathrm{smooth}}$ bounding smooth-min-entropy estimation error~\cite{lim2014,tomamichel2012}. The CHSH and QBER abort thresholds enter the budget through $\varepsilon_{\mathrm{PE}}$. For a verification subset of size $n_v$, the probability that an empirical statistic falls on the ``pass'' side of the threshold while the true value falls on the ``fail'' side decays exponentially in $n_v \cdot (\mathrm{threshold} - \mathrm{true})^2$ via Hoeffding-type tail inequalities. Under a four-way uniform budget allocation, $\varepsilon_{\mathrm{PE}} \leq \varepsilon_{\mathrm{sec}}/4 \approx 2.5 \times 10^{-11}$. Default thresholds (CHSH=$2.40$, QBER=$2.5\%$) sit well above the no-Eve operating point and satisfy this bound comfortably; optimised thresholds tighten the margin but remain within budget by construction.

\textbf{ML-DSA overhead.} Per protocol round ($\approx 1$\,s of quantum transmission):
\begin{equation}
t_{\mathrm{MLDSA}} = N_{\mathrm{exch}} \cdot (t_{\mathrm{sign}} + t_{\mathrm{verify}}) = 4 \cdot (0.23 + 0.19)~\mathrm{ms} = 1.68~\mathrm{ms} \label{eq:mldsa}
\end{equation}
i.e., $0.17$--$0.21\%$ of the round time across all distances evaluated.

All means use 95\% confidence intervals via Student's $t$-distribution over $N = 1{,}000$ trials. We use symmetric basis selection ($P = 1/3$ per basis). The intercept-resend attack is the standard worst-case individual-attack model; hardware-level attacks~\cite{lydersen2010,makarov2009} are discussed qualitatively in Section~\ref{sec:limitations}.

\subsection{Parameters}
Table~\ref{tab:params} lists all simulation parameters. The 100\,MHz pair rate reflects state-of-the-art SPDC sources deployed in early telecom-fiber E91 systems~\cite{poppe2004}.

\begin{table}[!t]
\caption{Simulation Parameters (Authors' Own Compilation)}
\label{tab:params}
\centering
\begin{tabular}{@{}lll@{}}
\toprule
\textbf{Parameter} & \textbf{Value} & \textbf{Notes} \\
\midrule
Entangled pair rate & 100~MHz & SPDC, telecom-fiber E91~\cite{poppe2004} \\
Fiber attenuation & 0.2~dB/km & Standard SMF-28 \\
Coupling efficiency & 10\% & Fiber coupling, polarisation optics \\
Detector efficiency & 85\% & InGaAs APD at 1{,}550~nm \\
Dark count probability & $10^{-5}$ & Per detection gate \\
Misalignment error & 1\% & Fixed optical misalignment \\
Basis selection prob. & 1/3 each & Symmetric 3-basis E91 \\
Eavesdropping fraction $p$ & 0, 0.05, 0.20 & Intercept-resend attack \\
ML-DSA sign / verify & 0.23 / 0.19~ms & NIST Level~3 reference impl. \\
CHSH abort threshold & $|S| < 2.40$ & Optimal: 2.66 (Sec.~\ref{sec:thresholds}) \\
QBER abort threshold & QBER $> 2.5\%$ & Optimal: 2.33\% (Sec.~\ref{sec:thresholds}) \\
Verification subset & 25\% of pulses & Used for CHSH test \\
LDPC efficiency $f$ & 1.16 & \cite{elkouss2011} \\
Privacy amplification & SHA-256 & Universal hash extractor \\
Security parameter $\varepsilon_{\mathrm{sec}}$ & $10^{-10}$ & Finite-key analysis \\
Trials per configuration & 1{,}000 & Independent; seed = 42 \\
\bottomrule
\end{tabular}
\end{table}

\subsection{Key Rate Results}
Table~\ref{tab:keyrates} reports final secret key rates with 95\% confidence intervals. Rows showing $0.00$\,kbps at $p > 0$ indicate that the protocol correctly detected and aborted due to eavesdropping (the secure behaviour). The Pure PQC entry is constant at 122\,kbps regardless of distance, reflecting fixed ML-KEM-768 handshake latency as defined in Section~\ref{sec:methodology}.

\begin{table}[!t]
\caption{Secret Key Rates (kbps) with 95\% Confidence Intervals (Authors' Own Simulation Results)}
\label{tab:keyrates}
\centering
\begin{tabular}{@{}lllll@{}}
\toprule
\textbf{Distance} & \textbf{p} & \textbf{Pure QKD} & \textbf{Pure PQC} & \textbf{Hybrid} \\
\midrule
10~km & 0.00 & $589.7 \pm 2.9$ & 122 & $588.7 \pm 2.9$ \\
10~km & 0.05 & $185.3 \pm 14.7$ & 122 & $185.0 \pm 14.7$ \\
10~km & 0.20 & $0.00 \pm 0.00$ & 122 & $0.00 \pm 0.00$ \\
50~km & 0.00 & $14.82 \pm 0.07$ & 122 & $14.80 \pm 0.07$ \\
50~km & 0.05 & $4.50 \pm 0.37$ & 122 & $4.49 \pm 0.37$ \\
50~km & 0.20 & $0.00 \pm 0.00$ & 122 & $0.00 \pm 0.00$ \\
100~km & 0.00 & $0.15 \pm 0.00$ & 122 & $0.15 \pm 0.00$ \\
100~km & 0.05 & $0.04 \pm 0.00$ & 122 & $0.04 \pm 0.00$ \\
100~km & 0.20 & $0.00 \pm 0.00$ & 122 & $0.00 \pm 0.00$ \\
\bottomrule
\end{tabular}
\end{table}

The hybrid's ML-DSA overhead is $0.17\%$ at 10\,km, $0.20\%$ at 50\,km, and $0.21\%$ at 100\,km---negligible in all cases.

\begin{figure}[!t]
\centering
\includegraphics[width=\columnwidth]{figures/fig3_keyrate.png}
\caption{Secret key rate comparison across distances and eavesdropping fractions. Each panel shows one distance (10~km, 50~km, 100~km); bars compare Pure QKD, Pure PQC (ML-KEM-768), and Hybrid (E91+ML-DSA) at $p = 0.00, 0.05, 0.20$. Error bars are 95\% CIs over 1{,}000 trials. (Authors' own work, generated from the simulator of~\cite{orthi_rane_zenodo}.)}
\label{fig:keyrate}
\end{figure}

\subsection{CHSH and QBER Detection Results}
\label{sec:detection}
Table~\ref{tab:detection} reports CHSH $|S|$, QBER, and per-trial detection probabilities for the hybrid at 10\,km.

\begin{table}[!t]
\caption{CHSH $|S|$, QBER, and Detection Probability vs. Eavesdropping Fraction (10~km) (Authors' Own Simulation Results)}
\label{tab:detection}
\centering
\small
\begin{tabular}{@{}llllll@{}}
\toprule
\textbf{p} & \textbf{Mean $|S|$} & \textbf{Mean QBER} & \textbf{$P_{\mathrm{CHSH}}$} & \textbf{$P_{\mathrm{QBER}}$} & \textbf{$P_{\mathrm{any}}$} \\
\midrule
0.00 & $2.800 \pm 0.002$ & $1.05\% \pm 0.028\%$ & 0.000 & 0.000 & 0.000 \\
0.05 & $2.729 \pm 0.002$ & $2.69\% \pm 0.032\%$ & 0.000 & 0.645 & 0.645 \\
0.20 & $2.521 \pm 0.002$ & $7.65\% \pm 0.031\%$ & 0.000 & 1.000 & 1.000 \\
\bottomrule
\end{tabular}
\end{table}

The results reveal complementary detection roles. At $p = 0.05$, $|S| = 2.729$ remains above the $2.40$ threshold; CHSH does not trigger. However, QBER ($2.69\%$) exceeds the $2.5\%$ threshold in $\sim 64.5\%$ of trials. At $p = 0.20$, QBER exceeds the threshold in every trial. Analytical solution of~\eqref{eq:S_of_p} gives $|S(p)| = 2.40$ at $p \approx 0.29$; thus for $p \geq 0.29$, CHSH provides additional confirmation at a different sensitivity range from QBER.

\begin{figure}[!t]
\centering
\includegraphics[width=\columnwidth]{figures/fig4_detection.png}
\caption{Eavesdropping detection performance. (a) Detection probability vs.\ eavesdropping fraction $p$ for QBER-only, CHSH-only, and combined abort policy, with simulation points from Table~\ref{tab:detection} overlaid. (b) CHSH $|S|$ (left axis) and QBER (right axis) vs.\ $p$ with shaded abort zones. QBER and CHSH have different operational sensitivity regions. Consistent colour palette throughout: QBER in amber, CHSH in navy, hybrid in green. Curves computed over $N = 1{,}000$ trials per point. (Authors' own work.)}
\label{fig:detection}
\end{figure}

\subsection{Finite-Key Analysis and Key Rate vs. Distance}
\label{sec:finitekey}
Table~\ref{tab:finitekey} reports finite-key results using the rigorous composable bound~\eqref{eq:finite_rig} with $\varepsilon_{\mathrm{sec}} = 10^{-10}$. The minimum block size $n_{\min} \approx 9.92 \times 10^3$~bits is distance-independent because the correction term depends only on block size. For comparison, the simplified $\sqrt{n}$ heuristic yields $n_{\min} \approx 1.12 \times 10^5$~bits---about an order of magnitude larger---confirming that the heuristic substantially over-estimates the finite-key penalty.

\begin{table}[!t]
\caption{Finite-Key Analysis ($p = 0$, $\varepsilon_{\mathrm{sec}} = 10^{-10}$) (Authors' Own Simulation Results)}
\label{tab:finitekey}
\centering
\begin{tabular}{@{}lllll@{}}
\toprule
\textbf{Distance} & \textbf{Asymptotic R} & \textbf{$n_{\min}$} & \textbf{Practical} & \textbf{Time to} \\
& \textbf{(kbps)} & \textbf{(bits)} & \textbf{block (Mb)} & \textbf{accumulate} \\
\midrule
10~km & 592.5 & $9.92 \times 10^3$ & 0.10 & $\sim 0.1$\,s \\
30~km & 93.9 & $9.92 \times 10^3$ & 0.10 & $\sim 0.7$\,s \\
50~km & 14.9 & $9.92 \times 10^3$ & 0.10 & $\sim 4.1$\,s \\
75~km & 1.49 & $9.92 \times 10^3$ & 0.10 & $\sim 41.2$\,s \\
100~km & 0.149 & $9.92 \times 10^3$ & 0.10 & $\sim 411.7$\,s \\
\bottomrule
\end{tabular}
\end{table}

\textbf{Practical operating range.} At 10--50\,km, the finite-key penalty is small relative to the asymptotic rate ($<\!2\%$ reduction) and the minimum block accumulates in seconds. At 100\,km, the bare $n_{\min}$ requires $\sim 6.9$\,minutes to accumulate and a practical $0.10$\,Mbit block requires comparable time; combined with the $0.149$\,kbps asymptotic rate, 100\,km remains at the edge of practicality. It is included in Table~\ref{tab:keyrates} for completeness; the recommended deployment range is $\leq 50$\,km.

\begin{figure}[!t]
\centering
\includegraphics[width=\columnwidth]{figures/fig5_distance.png}
\caption{Secret key rate vs.\ distance ($p = 0$, asymptotic regime). The hybrid (green) tracks Pure QKD (navy) closely, crossing below the PQC-only floor (amber dashed, 122~kbps) at $d \approx 27.4$~km. Simulation points from Table~\ref{tab:keyrates} are overlaid (95\% CIs over 1{,}000 trials). Colour palette consistent with Fig.~\ref{fig:detection}: hybrid green, QKD navy, PQC amber. The shaded region marks the practical metropolitan deployment range ($\leq 50$\,km). (Authors' own work.)}
\label{fig:distance}
\end{figure}

\subsection{Abort Threshold Co-Optimisation}
\label{sec:thresholds}
Fig.~\ref{fig:threshold} presents a 2D grid search over CHSH threshold $\in [2.20, 2.70]$ and QBER threshold $\in [1.5\%, 4.0\%]$, evaluating detection probability at $p = 0.05$ and false-positive (FP) rate at $p = 0$ across 625 threshold pairs.

The optimised thresholds are CHSH=$2.66$, QBER=$2.33\%$. Table~\ref{tab:thresholdcompare} compares default and optimised operating points across two methodologies and includes the empirical 100{,}000-trial validation result.

\begin{table}[!t]
\caption{Default vs.\ Optimised Operating Point Comparison (Authors' Own Simulation Results)}
\label{tab:thresholdcompare}
\centering
\small
\begin{tabular}{@{}p{3.1cm}p{2.0cm}p{2.0cm}@{}}
\toprule
\textbf{Metric} & \textbf{Default} & \textbf{Optimised} \\
& \textbf{(2.40, 2.5\%)} & \textbf{(2.66, 2.33\%)} \\
\midrule
Empirical FP, $p=0$, $N=1{,}000$ & $0/1000 = 0\%$ (Wilson upper $\approx 0.37\%$) & see F2-bis row \\
\midrule
Empirical FP, $p=0$, $N=10^5$ (F2-bis) & not separately measured & $392/100{,}000 = 0.392\%$ (95\% Wilson CI $[0.355\%, 0.433\%]$) \\
\midrule
Analytical FP (Gaussian tail) & $\sim 0.06\%$ & $0.40\%$ (within empirical CI) \\
\midrule
$P(\mathrm{detect} | p = 0.05)$ & $64.5\%$ & $75.6\%$ ($+11.1$ pp / $17.2\%$ rel.) \\
\midrule
$P(\mathrm{detect} | p = 0.20)$ & $100\%$ & $100\%$ (unchanged) \\
\bottomrule
\end{tabular}
\end{table}

\textbf{Reconciliation of false-positive estimates.} At the conservative defaults, the empirical $0/1{,}000$ result and the analytical $\sim 0.06\%$ Gaussian-tail estimate are mutually consistent (a true FP rate of $\sim 0.06\%$ produces $\sim 0.6$ events in 1{,}000 trials, within Poisson variance of zero observations). At the optimised thresholds, an expanded empirical validation at $N = 10^5$ trials (same Python implementation and random seed family, $p = 0$) produced 392 false aborts, yielding an empirical FP of $0.392\%$ with 95\% Wilson CI $[0.355\%, 0.433\%]$. The analytical prediction of $0.40\%$ falls inside this CI, confirming the Gaussian-tail model as a reliable predictor at this operating point. Disaggregating: 376 trials triggered QBER alone ($0.376\%$), 16 trials triggered CHSH alone ($0.016\%$), and no trial triggered both---matching the expected behaviour where QBER is the primary trigger and CHSH-only contributes a $\sim 3.5\sigma$ Gaussian tail (theoretical $0.023\%$; empirical $0.016\%$).

\textbf{Detection-probability improvement.} The optimised operating point improves $P(\mathrm{detect} | p = 0.05)$ from $64.5\%$ to $75.6\%$: $+11.1$ percentage points absolute ($75.6 - 64.5 = 11.1$), equivalently $17.2\%$ relative improvement. The exchange of $0.40$ percentage points of FP for $11.1$ percentage points of detection gain is a $\sim 28\!:\!1$ sensitivity/specificity trade.

\begin{figure}[!t]
\centering
\includegraphics[width=\columnwidth]{figures/fig6_threshold.png}
\caption{Abort threshold co-optimisation. (a) Detection probability at $p = 0.05$ as a function of CHSH and QBER thresholds; red contour marks the FP $= 0.5\%$ boundary; $\bigstar$ marks the optimal operating point (CHSH = 2.66, QBER = 2.33\%). (b) False-positive abort rate across the same grid. (Authors' own work.)}
\label{fig:threshold}
\end{figure}

\subsection{Quantum Circuit Validation via Qiskit Statevector Simulation}
To provide circuit-level validation of the analytical CHSH model~\eqref{eq:S_of_p}, we implemented the E91 Bell-state measurement protocol as a quantum circuit using Qiskit~2.4~\cite{qiskit} and executed it on the Qiskit Aer statevector simulator. The circuit prepares the $|\Phi^+\rangle = (|00\rangle + |11\rangle)/\sqrt{2}$ Bell state and measures Alice and Bob at the standard CHSH basis angles, computing $E(a,b)$ from 8{,}192 measurement shots per basis pair. Eve's intercept-resend attack is modelled in Qiskit as a partial $R_y$ rotation on Bob's qubit proportional to $p_{\mathrm{eve}}$, approximating the depolarising effect of full random-basis intercept-resend. Table~\ref{tab:qiskit} presents the resulting CHSH $|S|$ values alongside~\eqref{eq:S_of_p}.

\begin{table}[!t]
\caption{Qiskit Statevector Validation of CHSH Analytical Model. Eve's attack is modelled as a partial $R_y$ rotation on Bob's qubit proportional to $p_{\mathrm{eve}}$, approximating (rather than exactly simulating) full random-basis intercept-resend; Qiskit values therefore over-estimate $|S|$ relative to the analytical model. Both cross the 2.40 abort threshold in the same $p \approx 0.29$ region. (Authors' own work.)}
\label{tab:qiskit}
\centering
\begin{tabular}{@{}llll@{}}
\toprule
\textbf{$p_{\mathrm{eve}}$} & \textbf{$|S|$ Qiskit} & \textbf{$|S|$ Analytical} & \textbf{Deviation} \\
\midrule
0.00 & 2.8215 & 2.8000 & 0.021 \\
0.05 & 2.8196 & 2.7307 & 0.089 \\
0.10 & 2.7886 & 2.6614 & 0.127 \\
0.15 & 2.7510 & 2.5921 & 0.159 \\
0.20 & 2.6797 & 2.5228 & 0.157 \\
0.25 & 2.6113 & 2.4535 & 0.158 \\
0.30 & 2.5261 & 2.3842 & 0.142 \\
\bottomrule
\end{tabular}
\end{table}

The Qiskit circuit and~\eqref{eq:S_of_p} agree within $\pm 0.159$ $|S|$ units across all tested fractions. Both cross below the 2.40 abort threshold in the same $p \geq 0.29$ region. The analytical model is the more conservative estimate for threshold design.

\begin{figure}[!t]
\centering
\includegraphics[width=\columnwidth]{figures/fig7_qiskit.png}
\caption{Qiskit statevector validation of the analytical CHSH model. Green circles show $|S|$ measured from E91 Bell-state quantum circuits (8{,}192 shots per basis pair). Navy squares show~\eqref{eq:S_of_p} at the same fractions. Both cross the 2.40 abort threshold (red dotted) in the same $p \approx 0.29$ region. Maximum deviation is $\pm 0.159$ $|S|$ units. (Authors' own work.)}
\label{fig:qiskit}
\end{figure}

% ============================================================
\section{Discussion}
\label{sec:discussion}
% ============================================================

\subsection{Hybrid vs.\ Pure PQC and Pure QKD}
Pure QKD without authentication is operationally insecure: a man-in-the-middle who intercepts only the classical channel can substitute basis announcements and extract the key without touching the quantum channel. Adding ML-DSA closes this attack at negligible cost ($< 0.25\%$ overhead). Pure PQC provides no eavesdropping detection: if a future cryptanalytic advance compromises the KEM, the attacker learns the key without observable anomaly. The hybrid mitigates this risk because CHSH and QBER monitoring continue to provide quantum-layer anomaly visibility even if the classical layer is compromised. Implementation-level attacks such as detector blinding~\cite{lydersen2010,makarov2009} still require hardware-level countermeasures.

\subsection{Complementarity of CHSH and QBER Detection}
A key finding is that CHSH and QBER play partially complementary operational roles. For $p \leq 0.20$, QBER is the primary mechanism; for $p \geq 0.29$, CHSH provides additional confirmation. Under the intercept-resend model, simultaneous maintenance of QBER $< 2.5\%$ and $|S| > 2.40$ was not observed for $p \geq 0.05$ across all 1{,}000 trials.

\subsection{When Is the Hybrid Worth the Cost?}
For routine applications, PQC alone after NIST standardisation is sufficient and far less expensive. QKD hardware requires single-photon detectors, entangled sources, and dedicated or dark fiber. The hybrid is best suited for high-assurance communications where data value justifies the hardware cost, harvest-now-decrypt-later threat models apply~\cite{mosca2018}, and real-time intrusion detection is required.

\subsection{Practical Deployment Range}
The practical operational range is up to $50$\,km, yielding $\sim 14.8$\,kbps with no eavesdropping---sufficient to refresh AES-256 session keys approximately every 15~seconds. 100\,km exceeds practical range as discussed in Section~\ref{sec:finitekey}. ETSI GR QKD~004~\cite{etsi2020} and NIST FIPS 203/204~\cite{fips203,fips204} provide a clear standardisation path; a hybrid profile for TLS~1.3 or IPsec is a natural target.

\subsection{Limitations}
\label{sec:limitations}
Our simulation assumes an ideal entangled source producing pure Bell states. Real sources produce mixed states with reduced fidelity, lowering the maximum achievable $|S|$. The treatment of realistic devices by Xu \emph{et~al.}~\cite{xu2020} catalogues principal imperfections (efficiency mismatch, dead time, after-pulsing, polarisation-mode dispersion) that should be incorporated in any hardware-realistic refinement. The hardware implementation of Garms \emph{et~al.}~\cite{garms2024} captures these where our discrete-event model does not. Our simulation models only intercept-resend; more sophisticated strategies---including detector blinding~\cite{lydersen2010}, photon-number splitting, and trojan-horse attacks---require hardware-level countermeasures. Measurement-device-independent QKD~\cite{lo2012} and twin-field QKD~\cite{lucamarini2018} offer alternative architectures inherently resistant to detector-side attacks. The security argument of Section~\ref{sec:architecture} is informal, not a formal UC proof; Portmann and Renner~\cite{portmann2022} and Lim \emph{et~al.}~\cite{lim2014} provide composable frameworks within which a rigorous proof should be developed.

We additionally note three modelling assumptions that bound the generality of our results. First, a \emph{finite-key composability gap} remains: although Section~\ref{sec:methodology} now reports the rigorous Tomamichel-style bound~\eqref{eq:finite_rig}, the end-to-end protocol is not yet proven $\varepsilon$-secure in a fully composable sense, since parameter estimation, reconciliation, and privacy amplification are budgeted by union bound rather than via a single composable reduction. Second, we assume \emph{symmetric basis choice} ($P = 1/3$ per basis); biased-basis schemes can increase the sifting efficiency and would shift the key-rate and finite-key results. Third, we assume \emph{independent and identically distributed (IID) channel statistics} across rounds; correlated noise, drift, or non-Markovian memory effects in real fiber would invalidate the per-trial Gaussian model and require de Finetti or entropy-accumulation treatments for rigorous bounds.

% ============================================================
\section{Conclusion}
\label{sec:conclusion}
% ============================================================
The two roads to quantum safety are not forks in the road but parallel lanes on the same highway. QKD provides physics-based, computational-assumption-free eavesdropping detection with finite statistical confidence bounded by the verification block size. ML-DSA and its PQC siblings provide computationally secure authentication deployable on existing infrastructure. Neither is complete without the other.

This paper makes four specific contributions: it establishes the KEM/signature design distinction; provides a simulation-validated reference protocol with full methodology; quantifies the different operational sensitivity regions of QBER and CHSH detectors; and presents finite-key analysis, a continuous key-rate curve, and threshold co-optimisation showing CHSH=$2.66$ and QBER=$2.33\%$ improve simulated detection at $p = 0.05$ from $64.5\%$ to $75.6\%$ ($+11.1$ pp absolute, $17.2\%$ relative). A 100{,}000-trial empirical validation confirms the analytical FP prediction within sampling precision. ML-DSA overhead is under $0.25\%$ per round. The hybrid is practical for metropolitan-area deployments up to $50$\,km.

\textbf{Future work.} F1: implement on a real QKD testbed under realistic attacks. F2: develop a formal UC security proof. F3: integrate HQC encryption alongside ML-DSA authentication. F4: develop a standards-track hybrid profile for TLS~1.3/IPsec. F5: adapt to CV-QKD systems. F6: investigate joint optimisation of CHSH/QBER abort thresholds using the analytical crossover model across all $p$.

% ============================================================
\appendices
\section{Derivation of the QBER Intercept-Resend Term}
\label{app:qber}
% ============================================================
The eavesdropping contribution $p_{\mathrm{eve}}/3$ to the quantum bit error rate in~\eqref{eq:qber} follows from the random-basis intercept-resend model for three-basis E91~\cite{fuchs1997,fung2006}. Consider the fraction $p_{\mathrm{eve}}$ of photons that Eve intercepts. For each, Eve measures in a basis chosen at random from the three available and resends a freshly prepared state in that basis. Two events must combine to produce an induced error:
\begin{enumerate}
\item \textbf{Eve picks the wrong basis.} With three mutually non-orthogonal measurement bases, the probability that Eve's randomly chosen basis differs from the one Alice and Bob later use for key extraction is $2/3$. When Eve happens to choose the correct basis, her measurement and resend do not disturb the outcome and introduce no error.
\item \textbf{A wrong-basis measurement yields a random bit.} When Eve measures in the wrong (conjugate) basis, the resent photon collapses to a state uncorrelated with the original; Bob's subsequent measurement therefore yields the wrong bit with probability $1/2$.
\end{enumerate}
Multiplying the intercept fraction by these two conditional probabilities gives the per-photon induced error
\begin{equation}
e_{\mathrm{eve}} = p_{\mathrm{eve}} \cdot \frac{2}{3} \cdot \frac{1}{2} = \frac{p_{\mathrm{eve}}}{3}, \label{eq:qber_deriv}
\end{equation}
which, added to the baseline misalignment and dark-count terms, yields the QBER model of~\eqref{eq:qber}. This $p_{\mathrm{eve}}/3$ scaling is the QBER counterpart of the CHSH degradation in~\eqref{eq:S_of_p}; both arise from the same disturbance mechanism and together produce the complementary detection behaviour characterised in Section~\ref{sec:detection}.

% ============================================================
\section*{Acknowledgment}
% ============================================================
The authors thank the quantum cryptography and post-quantum cryptography communities for their open research, standardisation work, and freely available reference implementations.

% ============================================================
\section*{Data and Code Availability}
% ============================================================
The simulation code, all generated data files, and output figures are publicly available at GitHub (\url{https://github.com/Orthimalek/quantum-safety-hybrid-qkd-pqc}) and Zenodo (\url{https://doi.org/10.5281/zenodo.20144465}). The simulation is fully reproducible using Python~3 with NumPy and SciPy (random seed 42).

% ============================================================
% REFERENCES
% ============================================================
\begin{thebibliography}{40}

\bibitem{bennett1984}
C. H. Bennett and G. Brassard, ``Quantum cryptography: Public key distribution and coin tossing,'' in \emph{Proc. IEEE Int. Conf. Comput. Syst. Signal Process. (ICCSSP)}, Bangalore, India, Dec. 1984, pp. 175--179.

\bibitem{ekert1991}
A. K. Ekert, ``Quantum cryptography based on Bell's theorem,'' \emph{Phys. Rev. Lett.}, vol. 67, no. 6, pp. 661--663, Aug. 1991, doi: \href{https://doi.org/10.1103/PhysRevLett.67.661}{10.1103/PhysRevLett.67.661}.

\bibitem{shor1997}
P. W. Shor, ``Polynomial-time algorithms for prime factorization and discrete logarithms on a quantum computer,'' \emph{SIAM J. Comput.}, vol. 26, no. 5, pp. 1484--1509, Oct. 1997, doi: \href{https://doi.org/10.1137/S0097539795293172}{10.1137/S0097539795293172}.

\bibitem{grover1996}
L. K. Grover, ``A fast quantum mechanical algorithm for database search,'' in \emph{Proc. 28th Annu. ACM Symp. Theory Comput. (STOC)}, Philadelphia, PA, May 1996, pp. 212--219, doi: \href{https://doi.org/10.1145/237814.237866}{10.1145/237814.237866}.

\bibitem{nist2022}
National Institute of Standards and Technology, ``Status report on the third round of the NIST post-quantum cryptography standardization process,'' \emph{NIST Interagency Rep.}~8413, Jul. 2022, doi: \href{https://doi.org/10.6028/NIST.IR.8413}{10.6028/NIST.IR.8413}.

\bibitem{fips204}
National Institute of Standards and Technology, ``Module-lattice-based digital signature standard,'' \emph{FIPS} 204, Aug. 2024, doi: \href{https://doi.org/10.6028/NIST.FIPS.204}{10.6028/NIST.FIPS.204}.

\bibitem{fips203}
National Institute of Standards and Technology, ``Module-lattice-based key-encapsulation mechanism standard,'' \emph{FIPS} 203, Aug. 2024, doi: \href{https://doi.org/10.6028/NIST.FIPS.203}{10.6028/NIST.FIPS.203}.

\bibitem{hqc2022}
C. Aguilar Melchor \emph{et al.}, ``HQC (Hamming Quasi-Cyclic),'' NIST Post-Quantum Cryptography Round~4 Submission, 2022. [Online]. Available: \url{https://pqc-hqc.org/}

\bibitem{lo2005}
H.-K. Lo, X. Ma, and K. Chen, ``Decoy state quantum key distribution,'' \emph{Phys. Rev. Lett.}, vol. 94, no. 23, p. 230504, Jun. 2005, doi: \href{https://doi.org/10.1103/PhysRevLett.94.230504}{10.1103/PhysRevLett.94.230504}.

\bibitem{scarani2009}
V. Scarani \emph{et al.}, ``The security of practical quantum key distribution,'' \emph{Rev. Mod. Phys.}, vol. 81, no. 3, pp. 1301--1350, Sep. 2009, doi: \href{https://doi.org/10.1103/RevModPhys.81.1301}{10.1103/RevModPhys.81.1301}.

\bibitem{lim2014}
C. C. W. Lim, M. Curty, N. Walenta, F. Xu, and H. Zbinden, ``Concise security bounds for practical decoy-state quantum key distribution,'' \emph{Phys. Rev. A}, vol. 89, no. 2, p. 022307, Feb. 2014, doi: \href{https://doi.org/10.1103/PhysRevA.89.022307}{10.1103/PhysRevA.89.022307}.

\bibitem{gisin2002}
N. Gisin, G. Ribordy, W. Tittel, and H. Zbinden, ``Quantum cryptography,'' \emph{Rev. Mod. Phys.}, vol. 74, no. 1, pp. 145--195, Jan. 2002, doi: \href{https://doi.org/10.1103/RevModPhys.74.145}{10.1103/RevModPhys.74.145}.

\bibitem{liao2017}
S.-K. Liao \emph{et al.}, ``Satellite-based entanglement distribution over 1200 kilometers,'' \emph{Science}, vol. 356, no. 6343, pp. 1140--1144, Jun. 2017, doi: \href{https://doi.org/10.1126/science.aan3211}{10.1126/science.aan3211}.

\bibitem{chsh1969}
J. F. Clauser, M. A. Horne, A. Shimony, and R. A. Holt, ``Proposed experiment to test local hidden-variable theories,'' \emph{Phys. Rev. Lett.}, vol. 23, no. 15, pp. 880--884, Oct. 1969, doi: \href{https://doi.org/10.1103/PhysRevLett.23.880}{10.1103/PhysRevLett.23.880}.

\bibitem{pirandola2020}
S. Pirandola \emph{et al.}, ``Advances in quantum cryptography,'' \emph{Adv. Opt. Photon.}, vol. 12, no. 4, pp. 1012--1236, Dec. 2020, doi: \href{https://doi.org/10.1364/AOP.361502}{10.1364/AOP.361502}.

\bibitem{xu2020}
F. Xu, X. Ma, Q. Zhang, H.-K. Lo, and J.-W. Pan, ``Secure quantum key distribution with realistic devices,'' \emph{Rev. Mod. Phys.}, vol. 92, no. 2, p. 025002, May 2020, doi: \href{https://doi.org/10.1103/RevModPhys.92.025002}{10.1103/RevModPhys.92.025002}.

\bibitem{mosca2018}
M. Mosca, ``Cybersecurity in an era with quantum computers: Will we be ready?'' \emph{IEEE Security Privacy}, vol. 16, no. 5, pp. 38--41, Sep./Oct. 2018, doi: \href{https://doi.org/10.1109/MSP.2018.3761723}{10.1109/MSP.2018.3761723}.

\bibitem{bindel2019}
N. Bindel, J. Brendel, M. Fischlin, and B. Warinschi, ``Hybrid key encapsulation mechanisms and authenticated key exchange,'' in \emph{Post-Quantum Cryptography (PQCrypto 2019)}, \emph{Lect. Notes Comput. Sci.}, vol. 11505, Cham, Switzerland: Springer, 2019, pp. 206--226, doi: \href{https://doi.org/10.1007/978-3-030-25510-7_12}{10.1007/978-3-030-25510-7\_12}.

\bibitem{stebila2010}
D. Stebila, M. Mosca, and N. L\"utkenhaus, ``The case for quantum key distribution,'' in \emph{Quantum Communication and Quantum Networking}, vol. 36, Berlin, Germany: Springer, 2010, pp. 283--296, doi: \href{https://doi.org/10.1007/978-3-642-11731-2_35}{10.1007/978-3-642-11731-2\_35}.

\bibitem{etsi2020}
European Telecommunications Standards Institute, ``Quantum key distribution (QKD): Application interface,'' ETSI Group Rep.\ GR QKD 004, V2.1.1, Aug. 2020.

\bibitem{elkouss2011}
D. Elkouss, J. Martinez-Mateo, and V. Martin, ``Information reconciliation for quantum key distribution,'' \emph{Quantum Inf. Comput.}, vol. 11, no. 3--4, pp. 226--238, Mar. 2011, doi: \href{https://doi.org/10.26421/QIC11.3-4-2}{10.26421/QIC11.3-4-2}.

\bibitem{treiber2009}
A. Treiber \emph{et al.}, ``A fully automated entanglement-based quantum cryptography system for telecom fiber networks,'' \emph{New J. Phys.}, vol. 11, p. 045013, Apr. 2009, doi: \href{https://doi.org/10.1088/1367-2630/11/4/045013}{10.1088/1367-2630/11/4/045013}.

\bibitem{hsieh2025}
Y.-C. Hsieh \emph{et al.}, ``Hybrid schemes of NIST post-quantum cryptography standard algorithms and quantum key distribution for key exchange and digital signature,'' \emph{arXiv}:2510.02379, Oct. 2025.

\bibitem{garms2024}
L. Garms \emph{et al.}, ``Experimental integration of quantum key distribution and post-quantum cryptography in a hybrid quantum-safe cryptosystem,'' \emph{Adv. Quantum Technol.}, vol. 7, no. 4, p. 2300304, Apr. 2024, doi: \href{https://doi.org/10.1002/qute.202300304}{10.1002/qute.202300304}.

\bibitem{zeng2024}
P. Zeng \emph{et al.}, ``Practical hybrid PQC-QKD protocols with enhanced security and performance,'' \emph{arXiv}:2411.01086, Nov. 2024.

\bibitem{poppe2004}
A. Poppe \emph{et al.}, ``Practical quantum key distribution with polarization entangled photons,'' \emph{Opt. Express}, vol. 12, no. 16, pp. 3865--3871, Aug. 2004, doi: \href{https://doi.org/10.1364/OPEX.12.003869}{10.1364/OPEX.12.003869}.

\bibitem{lydersen2010}
L. Lydersen \emph{et al.}, ``Hacking commercial quantum cryptography systems by tailored bright illumination,'' \emph{Nature Photon.}, vol. 4, pp. 686--689, Aug. 2010, doi: \href{https://doi.org/10.1038/nphoton.2010.214}{10.1038/nphoton.2010.214}.

\bibitem{lo2012}
H.-K. Lo, M. Curty, and B. Qi, ``Measurement-device-independent quantum key distribution,'' \emph{Phys. Rev. Lett.}, vol. 108, no. 13, p. 130503, Mar. 2012, doi: \href{https://doi.org/10.1103/PhysRevLett.108.130503}{10.1103/PhysRevLett.108.130503}.

\bibitem{lucamarini2018}
M. Lucamarini \emph{et al.}, ``Overcoming the rate--distance limit of quantum key distribution without quantum repeaters,'' \emph{Nature}, vol. 557, pp. 400--403, May 2018, doi: \href{https://doi.org/10.1038/s41586-018-0066-6}{10.1038/s41586-018-0066-6}.

\bibitem{portmann2022}
C. Portmann and R. Renner, ``Security in quantum cryptography,'' \emph{Rev. Mod. Phys.}, vol. 94, no. 2, p. 025008, Jun. 2022, doi: \href{https://doi.org/10.1103/RevModPhys.94.025008}{10.1103/RevModPhys.94.025008}.

\bibitem{orthi_rane_zenodo}
V. D. Rane and S. M. Orthi, ``Simulation code: Two roads to quantum safety---hybrid E91 QKD + ML-DSA,'' Zenodo, Jun. 2026, doi: \href{https://doi.org/10.5281/zenodo.20144465}{10.5281/zenodo.20144465}.

\bibitem{tomamichel2012}
M. Tomamichel, C. C. W. Lim, N. Gisin, and R. Renner, ``Tight finite-key analysis for quantum cryptography,'' \emph{Nature Commun.}, vol. 3, p. 634, Jan. 2012, doi: \href{https://doi.org/10.1038/ncomms1631}{10.1038/ncomms1631}.

\bibitem{tomamichel2017}
M. Tomamichel, J. Martinez-Mateo, C. Pacher, and D. Elkouss, ``Fundamental finite key limits for one-way information reconciliation in quantum key distribution,'' \emph{Quantum Inf. Process.}, vol. 16, no. 11, p. 280, Nov. 2017, doi: \href{https://doi.org/10.1007/s11128-017-1709-5}{10.1007/s11128-017-1709-5}.

\bibitem{makarov2009}
V. Makarov, ``Controlling passively quenched single photon detectors by bright light,'' \emph{New J. Phys.}, vol. 11, p. 065003, Jun. 2009, doi: \href{https://doi.org/10.1088/1367-2630/11/6/065003}{10.1088/1367-2630/11/6/065003}.

\bibitem{qiskit}
IBM Quantum, ``Qiskit: An open-source framework for quantum computing,'' 2024, doi: \href{https://doi.org/10.5281/zenodo.2573505}{10.5281/zenodo.2573505}.

\bibitem{fuchs1997}
C. A. Fuchs, N. Gisin, R. B. Griffiths, C.-S. Niu, and A. Peres, ``Optimal eavesdropping in quantum cryptography. I. Information bound and optimal strategy,'' \emph{Phys. Rev. A}, vol. 56, no. 2, pp. 1163--1172, Aug. 1997, doi: \href{https://doi.org/10.1103/PhysRevA.56.1163}{10.1103/PhysRevA.56.1163}.

\bibitem{devetak2005}
I. Devetak and A. Winter, ``Distillation of secret key and entanglement from quantum states,'' \emph{Proc. R. Soc. A}, vol. 461, no. 2053, pp. 207--235, Jan. 2005, doi: \href{https://doi.org/10.1098/rspa.2004.1372}{10.1098/rspa.2004.1372}.

\bibitem{renner2008}
R. Renner, ``Security of quantum key distribution,'' \emph{Int. J. Quantum Inf.}, vol. 6, no. 1, pp. 1--127, Feb. 2008, doi: \href{https://doi.org/10.1142/S0219749908003256}{10.1142/S0219749908003256}.

\bibitem{fung2006}
C.-H. F. Fung, X. Ma, and H. F. Chau, ``Practical issues in quantum-key-distribution post-processing,'' \emph{Phys. Rev. A}, vol. 81, no. 1, p. 012318, Jan. 2010, doi: \href{https://doi.org/10.1103/PhysRevA.81.012318}{10.1103/PhysRevA.81.012318}.

\bibitem{hayashi2012}
M. Hayashi and T. Tsurumaru, ``Concise and tight security analysis of the Bennett--Brassard 1984 protocol with finite key lengths,'' \emph{New J. Phys.}, vol. 14, p. 093014, Sep. 2012, doi: \href{https://doi.org/10.1088/1367-2630/14/9/093014}{10.1088/1367-2630/14/9/093014}.

\end{thebibliography}

\end{document}
